% !TeX root = ../main.tex
% OOPSLA requires this section just before the references; it does not count towards the page
% limit. It must say whether an artifact exists, its nature and limitations, and whether it will
% be submitted for Artifact Evaluation. Anonymous builds name the supplementary material, which
% artifact/build.py writes; the camera-ready build names \artifactlink.
\section*{Data-Availability Statement}

\ifleananonymous
  An anonymized copy of the artifact is provided as supplementary material.
\else
  The artifact is available at \artifactlink.
\fi
It contains the Lean formalization and the command-line tool, the benchmarks, scripts and
recorded results of \cref{sec:evaluation}, the test suites, and the simulator. A website included
in it shows the simulator, the API documentation of the formalization and all source files, and
opens in a browser without installing anything.
Running \texttt{lake build -{}-wfail} checks every proof, and checks that each theorem in
\texttt{Theorems.lean}, which states the theorems of this paper, uses only Lean's standard axioms
(\texttt{propext}, \texttt{Classical.choice}, \texttt{Quot.sound}). The first build takes about
half an hour on a 16-core machine.
The Monte Carlo experiments use the compiled Lean runtime; plotting
requires Python and Matplotlib. Exact analysis using the Storm backend
also requires Storm. Timings depend on the machine.
